% =====================================================================================================
\section{Setting, the method, and the shared assumptions}\label{sec:setting}

This section fixes the data, the model, the posterior and the few objects of the method that the results talk about.
The methodology section describes the method in full; Appendix~\ref{app:notation} lists the symbols.

\paragraph{Notation.} $\set{k}:=\{1,\dots,k\}$; $\norm{\cdot}$ is the Euclidean norm; $\lmin(\cdot)$ is the smallest
eigenvalue of a symmetric matrix; $\abs{S}$ is the size of a set; $\ind{E}$ is an indicator; $a\vee b:=\max(a,b)$. A star marks a true value,
and the superscript $\mathrm r$ marks raw data. $V_n=\OP(a_n)$ means that $\abs{V_n}/a_n$ is bounded in probability, and $V_n=\oP(a_n)$ that it tends to zero in
probability; \emph{with probability tending to one} means that the probability of the event tends to one. For positive
sequences, $a_n=O(b_n)$ (or $a_n\lesssim b_n$, $b_n\gtrsim a_n$) means $a_n\le Cb_n$ for a constant $C$, $a_n=o(b_n)$
means $a_n/b_n\to0$, and $a_n\gg b_n$ means $b_n=o(a_n)$. A quantity is \emph{fixed} if
it does not depend on $n$; all results are asymptotic in $n\to\infty$, along a sequence of problems in which $p$, the
design and the coefficients may change with $n$.

\subsection{Data, model and posterior}\label{sec:model}

We observe a response $\yr\in\R^n$ and $\praw$ predictors. The method removes every predictor that is constant in the
sample (its effect is absorbed by the intercept); $\Xr$ is the $n\times p$ matrix of the $p$ retained predictors. The data
follow a Gaussian linear model with an intercept,
\begin{equation}\label{eq:model}
  \yr=\alpha^*\one_n+\Xr\bstar+\varepsilon,\qquad \varepsilon\sim\Normal_n(0,\sigs^2I_n),
\end{equation}
with $\sigs^2>0$ and $\varepsilon$ independent of $\Xr$. A \emph{model} is a set $S\subseteq\set{p}$; the \emph{true
model} is $\Sz:=\{j:\beta^*_j\neq0\}$, of size $\sz$. Its members are the \emph{true predictors}; the others are
\emph{null predictors}. In a \emph{fixed design} $\Xr$ is deterministic; in a \emph{random Gaussian design} its rows are
i.i.d.\ $\Normal_p(m_n,\Sig)$, and $W=(W_1,\dots,W_p)\sim\Normal_p(0,\Sig)$ denotes a centred copy of one row, with
sub-vectors $W_S$ and sub-matrices $\Sigma_{n,S}$.

\paragraph{What the method sees.} The method centres and scales every predictor and the response: $X:=P^\perp_{\one}\Xr
D^{-1}$ and $y:=P^\perp_{\one}\yr/s_y$, with $P^\perp_{\one}$ the centring matrix, $D$ the diagonal matrix of column sample
standard deviations and $s_y$ that of the response. Then $P^\perp_{\one}\yr=X\tstar+P^\perp_{\one}\varepsilon$ with the
\emph{standardised coefficients} $\tstar:=D\bstar$ (same support $\Sz$), the \emph{mean signal} $\mu:=X\tstar$, and
$\tmin:=\min_{j\in\Sz}\abs{\theta^*_j}$. For a model $S$, $X_S$ holds its columns, $G_S:=X_S^\top X_S$, $\Rhat_S:=G_S/(n-1)$
is its sample correlation matrix and $P_S$ the projection onto its column span.

\paragraph{The prior and the posterior, as coded.} The method fits a working Gaussian linear model,
$y\mid S,\theta_S,\sigma^2\sim\Normal_n(X_S\theta_S,\sigma^2I_n)$, to the standardised data; the theorems are proved
under the true model~\eqref{eq:model}, not under this working model. It considers the models of at most $K:=\min(p,\lfloor
n/2\rfloor)$ predictors, $\Om:=\{S:\abs{S}\le K\}$, and puts on them the beta-binomial prior $\BB(1,p)$ \citep{ScottBerger2010},
restricted to $\Om$; given $S$, a $g$-prior with $g=n$ and a unit ridge on the coefficients \citep{Zellner1986}; and an
inverse-gamma prior $\IG(0.01,0.01)$ on the noise variance. Integrating out the coefficients and the variance gives each
model a marginal likelihood $m(S)$ in closed form (methodology section, eq.~(6)), and the \emph{posterior over models}
\begin{equation}\label{eq:posterior}
  \pis{S}=\frac{e^{F(S)}}{Z},\qquad F(S):=\log\{\text{prior of }S\}+\log m(S),\qquad Z:=\sum_{S'\in\Om}e^{F(S')} .
\end{equation}
$F(S)$ is the \emph{log score} of $S$, and the normalising constant $Z$ is the \emph{evidence}. Adding one predictor to a model of size $k$ multiplies the prior by
$(k+1)/(2p-1-k)$, and, in a well-conditioned model, multiplies $m(S)$ by about $(n+1)^{-1/2}$ (an Occam factor) times a
gain of fit.

\paragraph{Posterior summaries.} The \emph{inclusion probability} of predictor $j$ is $\PIP_j:=\sum_{S\ni j}\pis{S}$;
the \emph{highest posterior model} is $\HPM:=\argmax_S\pis{S}$; the \emph{median probability model} is
$\MPM:=\{j:\PIP_j>1/2\}$ \citep{BarbieriBerger2004}; and the \emph{posterior miss} is $1-\pis{\Sz}$.

\begin{lemma}[The posterior miss controls the summaries]\label{lem:summaries}
For every data set, $\max_{j\in\set{p}}\abs{\PIP_j-\ind{j\in\Sz}}\le1-\pis{\Sz}$. If $\pis{\Sz}>1/2$, then
$\HPM=\MPM=\Sz$.
\end{lemma}

The proof is two lines: every model that misses a true predictor, or contains a null one, differs from $\Sz$.

\subsection{A run of the method, and what it returns}\label{sec:method}

A \emph{run} is one call of the software on a data set, with a seed and a computing allocation; a data set is
\emph{valid} when $n\ge3$, at least one predictor is given and every entry is finite. When $\Om$ is small enough the
method enumerates it; otherwise it runs the sampler. A run either \emph{produces output}, with status ``exact''
(enumeration) or ``completed'' (sampler; a \emph{completed run}), or stops without estimates. Four objects of a run appear
in the results.
\begin{itemize}[leftmargin=1.5em,itemsep=1pt]
  \item \emph{The start law.} The method partitions the predictors into blocks of at most $16$, by clustering on sample
    correlations, and scores every subset of each block exactly. The \emph{start law} is the posterior with all
    interactions between blocks $C_1,\dots,C_J$ ignored,
    $\pi_0(S)\propto w(\abs{S})\prod_{b=1}^{J}m(S\cap C_b)/m(\emptyset)$, where $w(\abs{S})$ is the prior weight of a
    model of that size; it can be sampled exactly. A maximiser of $\pi_0$ is a \emph{start mode} $\tSz$.
  \item \emph{The sampler.} A population of models moves from $\pi_0$ to the posterior through intermediate
    distributions, with temperatures and effort chosen by an internal \emph{calibration} that is frozen before the final
    run. The calibration \emph{succeeds} when it produces a sequence of temperatures (by its main construction or a
    fallback); otherwise the run stops. We write $\Acal$ for the event that the method uses the sampler and its
    calibration succeeds. The final run consists of four independent copies, the \emph{islands}.
  \item \emph{The reported model.} After the sampler, two local searches run (from the best model found, and from the
    start mode); each moves to the best neighbouring model (one predictor added, deleted or swapped) while that strictly
    improves the log score, for at most a capped number of moves ($\ell_n:=\lceil\log_2(n+p+1)\rceil$ for the first search,
    $\ell_n+\abs{\tSz}$ for the search from the start mode), and then scans the neighbours once more.
    The \emph{reported model} $\Shat$ is the best-scoring model the run has recorded. The method \emph{certifies}
    $\Shat$ (sets its local-optimality flag) when it enumerated $\Om$, or when $\Shat$ ends a local search whose last scan
    found no strict improvement; a search that reaches its cap and still finds an improving neighbour stops uncertified.
  \item \emph{The estimates.} Island $r$ estimates the inclusion probabilities and the evidence; the method pools the
    four islands with equal weights into $\PIPhat_j$ and $\Zhat$, and reports $\MPMhat:=\{j:\PIPhat_j>1/2\}$.
\end{itemize}

\subsection{Shared assumptions and conventions}\label{sec:standing}

The assumptions below are used across Sections~\ref{sec:model-result}--\ref{sec:estimates}; the others are stated
where they are first used. Our simulation designs (S1--S6, the strong-factor study, the designs IND, AR and BLK of the
planned runs at $p=10^6$, and FAC, a contingent design planned at $p=10^4$ only) are described in Assumptions~\S2: S1 and IND have independent predictors; S2--S4,
S6, AR (autoregressive) and BLK (blocks) are locally correlated; S5, the strong-factor study and FAC have pervasive common
factors.

Two words recur in Results~3 and~4. A true predictor is \emph{visible} if its signal does not cancel against any $15$ other
predictors (the method's blocks hold at most $16$); a null predictor \emph{leaks} if, given up to $15$ others, it can imitate
part of a true predictor's signal. Sections~\ref{sec:reported} and~\ref{sec:estimates} make both precise.
The assumptions section judges each of them and compares it with \citet{SongLiang2015}.

\begin{assumption}{GL}{Gaussian linear model}
The data follow~\eqref{eq:model}, with a fixed noise variance $\sigs^2>0$ and noise independent of the design. \hfill [\SA{3.1}]
\end{assumption}

\begin{assumption}{DIM}{Dimension}
$2\le p\le n^A$ for a fixed $A\ge1$. \hfill [\SA{3.1}]
\end{assumption}

\begin{assumption}{SP}{Sparsity}
The true model size $\sz\ge1$ is fixed. \hfill [\SA{3.1}]
\end{assumption}

\begin{assumption}{GR}{Gaussian rows}
The rows of $\Xr$ are i.i.d.\ $\Normal_p(m_n,\Sig)$, where $\Sig$ is a correlation matrix. \hfill [\SA{3.2}]
\end{assumption}

\begin{assumption}{EFS}{Population eigenvalue floor}
$\lmin(\Sigma_{n,S})\ge\varphi^2$ for every $S$ with $\abs{S}\le5(A+1)(\sz+1)\log n$, for a fixed $\varphi\in(0,1]$ (for
instance, $\lmin(\Sig)\ge\varphi^2$). \hfill [\SA{3.2}]
\end{assumption}

The conventions fix what ``the method'' means and idealise the computer; they are not statistical assumptions.
\emph{Convention \as{C1}} (the method): exact arithmetic, the default configuration, valid inputs and seed; a run that
meets it is \emph{admissible}, a sample size $n$ is admissible if some data set of that size is valid, and statements
hold for every valid seed and every outcome of the calibration.
\phantomsection\label{asm:C1}%
\emph{Convention \as{C2}} (delivery): an idealised machine that never overflows or fails, ideal random numbers, and an
allocation of memory and time no smaller than the method's own polynomial requirement formulas ask for; \emph{\as{C2p}} also lifts two fixed integer limits of the code (Appendix~\ref{app:conventions}).\phantomsection\label{asm:C2}\phantomsection\label{asm:C2p}
\emph{Convention \as{C3}} (sampler): exact arithmetic, ideal random numbers independent of the data and the calibration,
an exact start sampler, and Markov moves that do what their definitions say.\phantomsection\label{asm:C3}
Appendix~\ref{app:conventions} gives the full lists.
